% Principal: GDZ PPN236018647, 1917, PDF native 14–15, printed 7–8.
\SourceLeaf{14}\NoSourceHead
\pdfbookmark[1]{Vorwort zu Eisensteins Gesammelten Abhandlungen}{eisenstein}
\HistoricTitle{VORWORT.}
\HistoricCitation{Gesammelte Abhandlungen von G. \textsc{Eisenstein}, Berlin 1848. S. 1, 2.}

Die zuerst in den verschiedenen Bänden von \textsc{Crelle}'s Journal für Mathematik erschienenen und hier gesammelten Aufsätze bewegen sich, theils in der Höheren Arithmetik, theils in der Theorie der über Logarithmen und Kreisgrössen hinaus liegenden transcendenten Functionen, theils in der Verknüpfung dieser beiden grossen Gebiete, die zu den schönsten und fruchtbarsten im ganzen Umfange der Mathematik gehören. Die Höhere Arithmetik bietet einen unerschöpflichen Reichthum an interessanten Wahrheiten dar, und zwar an solchen, die nicht vereinzelt, sondern in innigem Zusammenhange stehen und immer neue, ja unerwartete Verknüpfungen erkennen lassen, je weiter die Wissenschaft sich ausbildet. Ein grosser Theil ihrer Lehren gewinnt auch einen neuen Reiz durch die Eigenthümlichkeit, dass gewichtige Lehrsätze in einfach ausgeprägtem Inhalt uns leicht durch Induction zugeführt werden, deren Begründung doch so tief liegt, dass man erst nach vielen vergeblichen Versuchen dazu gelangt, und dann meistens erst auf beschwerlichen künstlichen Wegen, während die einfacheren Methoden lange verborgen bleiben. Auch auf dem Felde der transcendenten Functionen fehlt es nicht an ähnlichen Reizen und ähnlichen Erscheinungen.

Von den eigenthümlichen Schönheiten dieser Gebiete haben Alle sich angezogen gefühlt, die darin beschäftigt gewesen sind: keiner hat es wohl so oft ausgesprochen wie \textsc{Euler}, der namentlich in fast allen seinen zahlreichen, zur Höheren Arithmetik gehörenden Aufsätzen die Erklärung wiederholt, wie viele Freude ihm diese Forschungen machen, und wie sehr er darin eine Erholung und eine Stärkung zu andern der unmittelbaren practischen Anwendung näher liegenden Arbeiten finde. Mit eben so grosser Lebhaftigkeit spricht er seine Überraschung aus, als zu seiner Kenntnis gekommen war,
% The paragraph continues across native 14–15 without a new historical start.
\par\clearpage\SourceLeaf{15}\SourceHead{8}\noindent
dass seine eigene Auflösung eines die transcendenten Functionen betreffenden Fundamen\-tal-Problems, mit welchem er sich viele Jahre beschäftigt hatte, an Einfachheit weit überboten sei durch eine neue Auflösung desselben Problems von \textsc{Lagrange}. »Penitus obstupui«, (sagt er Acta Acad. Petrop. T. 3.) »quum hoc mihi nunciaretur.«

Die vorliegenden Aufsätze enthalten so viel treffliches und gediegenes, dass durch dieselben dem Verfasser ein ehrenvoller Platz neben seinen Vorgängern gesichert wird, an deren Arbeiten jene sich würdig anschliessen. Ihre Zusammenstellung verpflichtet alle die zum Danke, denen sie dadurch zugänglicher werden.

\bigskip\noindent Göttingen im September 1847.
\begin{flushright}C. F. \textsc{Gauss}.\end{flushright}
\ClosingRule
\clearpage
